\begin{table}[htbp]\centering
\caption{H2. Relative choice: standardized repression minus standardized redistribution}
\label{tab:h2choice}
\small
\resizebox{\textwidth}{!}{%
\begin{tabular}{lcccc}
\toprule
 & (1) & (2) & (3) & (4) \\
\midrule
log K/L (t-1) & -0.203 (0.307) & 0.003 (0.178) & -0.059 (0.096) & -0.096 (0.355) \\
\quad wild-cluster bootstrap p & [0.541] & [0.984] & [0.597] & [0.828] \\
Threat: past mobilization (t-1) &  & -0.999** (0.453) & -0.665 (0.504) &  \\
K/L x past mobilization &  & 1.119*** (0.433) & 0.781 (0.501) &  \\
\quad wild-cluster bootstrap p &  & [0.048] & [0.246] &  \\
Threat: anti-system movements &  &  &  & 2.362 (1.922) \\
K/L x anti-system &  &  &  & -2.215 (2.103) \\
\quad wild-cluster bootstrap p &  &  &  & [0.415] \\
Military expenditure \% GDP (t-1) &  &  & -0.141*** (0.040) & -0.042 (0.116) \\
Military support (t-1) &  &  &  & 0.616 (0.395) \\
Driscoll-Kraay SE, log K/L &  &  &  & (0.316) \\
Observations & 553 & 285 & 232 & 493 \\
Countries & 18 & 11 & 10 & 17 \\
\bottomrule
\end{tabular}}
\begin{tablenotes}\footnotesize Dependent variable: z(physical repression) - z(absolute redistribution). Positive coefficients mean the policy mix tilts toward repression. Two-way FE, country-clustered SEs; Driscoll-Kraay SE (3 lags) reported for column (4). Columns (2)-(3) use the mobilization-based threat proxy (11 countries), column (4) the anti-system-movement proxy (18 countries).\end{tablenotes}
\end{table}